% =============================================================================
% preamble/floats.tex: figures, tables and algorithms ("floats": LaTeX moves
% them to where they fit best, so refer to them by number, not "below").
%
% Numbering is per chapter: Figure 3.1, Table 3.1, Algorithm 3.1. The book
% class does this for figures and tables; the [chapter] option below does it
% for algorithms.
%
% Conventions used here (and in most CS dissertations):
%   - table captions go ABOVE the table (put \caption before \begin{tabular})
%   - figure captions go BELOW the figure (put \caption after \includegraphics)
%   - put \label straight after \caption, so the label gets the right number
% Example code for all three is in README.md.
% =============================================================================

% --- Captions ------------------------------------------------------------------------------
\usepackage{caption}
\captionsetup{
  labelfont = bf,       % "Figure 3.1" in bold
  labelsep  = colon,    % "Figure 3.1: Caption text"
  font      = small,    % caption text one size smaller than body text
}
\captionsetup[table]{position=top}  % tells caption that table captions are
                                    % above, so it puts the gap below them

% --- Tables ---------------------------------------------------------------------------------
\usepackage{booktabs}   % professional horizontal rules: \toprule, \midrule,
                        % \bottomrule (avoid vertical lines in tables)
\usepackage{tabularx}   % tabularx environment: table of a set total width with
                        % X columns that wrap text and share the spare space
\usepackage{array}      % extra column options, e.g. p{3cm} with
                        % >{\raggedright\arraybackslash}

% --- Algorithms -----------------------------------------------------------------------------
% Two packages work together:
%   algorithm      the floating box with a caption and number, and the
%                  \listofalgorithms command used in frontmatter/lists.tex
%   algpseudocode  the pseudocode itself inside \begin{algorithmic}:
%                  \Require, \Ensure, \State, \If{..}..\EndIf, \For{..}..\EndFor,
%                  \While, \Function{Name}{args}..\EndFunction, \Return, \Comment
% \begin{algorithmic}[1] numbers every line; \label a line to \cref it.
\usepackage[chapter]{algorithm}  % [chapter]: number algorithms per chapter
\usepackage{algpseudocode}
